%% ============================================================
%%  SECCIÓN 7 — Anexos y Demostraciones Auxiliares
%% ============================================================

\begin{appendices}

\section{Demostración Matricial de la Identidad de Levi-Civita}
\label{anexo:levi_civita}

En esta sección demostraremos de forma elegante la identidad tensorial fundamental:
\begin{equation}
    \varepsilon_{ijk}\varepsilon_{lmk} = \delta_{il}\delta_{jm} - \delta_{im}\delta_{jl},
\end{equation}
sin recurrir a verificar manualmente los $3^4 = 81$ casos combinatorios \cite{sakurai}.

Consideremos una base ortonormal estándar de $\mathbb{R}^3$, $\{\hat{e}_1, \hat{e}_2, \hat{e}_3\}$, donde $\hat{e}_a \cdot \hat{e}_b = \delta_{ab}$. El pseudotensor de Levi-Civita puede definirse geométricamente como el determinante formado por tres vectores base:
\begin{equation}
    \varepsilon_{ijk} = \det \begin{pmatrix} \hat{e}_i & \hat{e}_j & \hat{e}_k \end{pmatrix}.
\end{equation}
Consideremos ahora el producto de dos símbolos de Levi-Civita con índices totalmente libres $(i, j, k)$ y $(l, m, n)$. Como el determinante del producto de dos matrices es el producto de sus determinantes ($\det(A B^T) = \det(A)\det(B)$):
\begin{equation}
    \varepsilon_{ijk}\varepsilon_{lmn} = \det \left[ \begin{pmatrix} \hat{e}_i^T \\ \hat{e}_j^T \\ \hat{e}_k^T \end{pmatrix} \begin{pmatrix} \hat{e}_l & \hat{e}_m & \hat{e}_n \end{pmatrix} \right] = \det \begin{pmatrix}
        \hat{e}_i \cdot \hat{e}_l & \hat{e}_i \cdot \hat{e}_m & \hat{e}_i \cdot \hat{e}_n \\
        \hat{e}_j \cdot \hat{e}_l & \hat{e}_j \cdot \hat{e}_m & \hat{e}_j \cdot \hat{e}_n \\
        \hat{e}_k \cdot \hat{e}_l & \hat{e}_k \cdot \hat{e}_m & \hat{e}_k \cdot \hat{e}_n
    \end{pmatrix}.
\end{equation}
Dado que $\hat{e}_a \cdot \hat{e}_b = \delta_{ab}$, esta matriz de $3 \times 3$ contiene exclusivamente deltas de Kronecker:
\begin{equation}
    \varepsilon_{ijk}\varepsilon_{lmn} = \begin{vmatrix}
        \delta_{il} & \delta_{im} & \delta_{in} \\
        \delta_{jl} & \delta_{jm} & \delta_{jn} \\
        \delta_{kl} & \delta_{km} & \delta_{kn}
    \end{vmatrix}.
    \label{eq:levi_civita_det3x3}
\end{equation}
Para obtener la contracción deseada, igualamos los índices $n = k$ y sumamos sobre $k \in \{1, 2, 3\}$:
\begin{equation}
    \sum_{k=1}^3 \varepsilon_{ijk}\varepsilon_{lmk} = \sum_{k=1}^3 \begin{vmatrix}
        \delta_{il} & \delta_{im} & \delta_{ik} \\
        \delta_{jl} & \delta_{jm} & \delta_{jk} \\
        \delta_{kl} & \delta_{km} & \delta_{kk}
    \end{vmatrix}.
\end{equation}
Desarrollando el determinante por la tercera columna mediante la regla de Laplace:
\begin{equation}
    \sum_{k=1}^3 \varepsilon_{ijk}\varepsilon_{lmk} = \sum_{k=1}^3 \left[ \delta_{ik} (\delta_{jl}\delta_{km} - \delta_{jm}\delta_{kl}) - \delta_{jk} (\delta_{il}\delta_{km} - \delta_{im}\delta_{kl}) + \delta_{kk} (\delta_{il}\delta_{jm} - \delta_{im}\delta_{jl}) \right].
\end{equation}
Evaluemos la suma sobre $k$ término a término aplicando las propiedades del delta de Kronecker:
\begin{enumerate}
    \item Para el primer término:
    \begin{equation}
        \sum_{k=1}^3 \delta_{ik} (\delta_{jl}\delta_{km} - \delta_{jm}\delta_{kl}) = \delta_{jl}\delta_{im} - \delta_{jm}\delta_{il}.
    \end{equation}
    \item Para el segundo término:
    \begin{equation}
        -\sum_{k=1}^3 \delta_{jk} (\delta_{il}\delta_{km} - \delta_{im}\delta_{kl}) = -(\delta_{il}\delta_{jm} - \delta_{im}\delta_{jl}) = -\delta_{il}\delta_{jm} + \delta_{im}\delta_{jl}.
    \end{equation}
    \item Para el tercer término, notemos que $\sum_{k=1}^3 \delta_{kk} = \delta_{11} + \delta_{22} + \delta_{33} = 3$:
    \begin{equation}
        \sum_{k=1}^3 \delta_{kk} (\delta_{il}\delta_{jm} - \delta_{im}\delta_{jl}) = 3 (\delta_{il}\delta_{jm} - \delta_{im}\delta_{jl}).
    \end{equation}
\end{enumerate}
Sumando algebraicamente las tres contribuciones:
\begin{align}
    \sum_{k=1}^3 \varepsilon_{ijk}\varepsilon_{lmk} &= (\delta_{jl}\delta_{im} - \delta_{jm}\delta_{il}) - (\delta_{il}\delta_{jm} - \delta_{im}\delta_{jl}) + 3(\delta_{il}\delta_{jm} - \delta_{im}\delta_{jl}) \nonumber\\
    &= -(\delta_{il}\delta_{jm} - \delta_{im}\delta_{jl}) - (\delta_{il}\delta_{jm} - \delta_{im}\delta_{jl}) + 3(\delta_{il}\delta_{jm} - \delta_{im}\delta_{jl}) \nonumber\\
    &= (-1 - 1 + 3)(\delta_{il}\delta_{jm} - \delta_{im}\delta_{jl}) \nonumber\\
    &= \boxed{\delta_{il}\delta_{jm} - \delta_{im}\delta_{jl}}.
\end{align}
La identidad queda formalmente probada de manera exacta.

\section{Formulario Maestro de Relaciones de Conmutación y Fórmulas Clave}
\label{anexo:formulario}

Para una rápida consulta previa a los exámenes y sesiones de resolución de problemas, la Tabla~\ref{tab:formulario_maestro} sintetiza las fórmulas fundamentales deducidas en este documento.

\begin{table}[H]
    \centering
    \caption{Relaciones operacionales y algebraicas fundamentales de la mecánica cuántica.}
    \label{tab:formulario_maestro}
    \vspace{0.2cm}
    \resizebox{\linewidth}{!}{%
    \begin{tabular}{@{}llp{8.5cm}@{}}
        \toprule
        \textbf{Categoría} & \textbf{Expresión Matemática} & \textbf{Significado Físico / Aplicación} \\
        \midrule
        \textbf{Postulado de Medida} & $\Delta\hat{A} \equiv \hat{A} - \ev{\hat{A}}\hat{I}$ & Operador desviación (fluctuación pura alrededor de la media). \\
        \addlinespace[2pt]
        & $(\Delta A)^2 = \ev{\hat{A}^2} - \ev{\hat{A}}^2$ & Varianza cuántica intrínseca del estado $\ket{\psi}$. \\
        \midrule
        \textbf{Incertidumbre} & $(\Delta A)^2 (\Delta B)^2 \ge \frac{1}{4}|\ev{\{\Delta\hat{A},\Delta\hat{B}\}}|^2 + \frac{1}{4}|\ev{[\hat{A},\hat{B}]}|^2$ & Relación generalizada de Robertson--Schrödinger (incluye covarianza). \\
        \addlinespace[2pt]
        & $\Delta A \, \Delta B \ge \frac{1}{2}|\ev{[\hat{A},\hat{B}]}|$ & Relación de Heisenberg (al despreciar covarianza cuántica). \\
        \midrule
        \textbf{Álgebra de Lie} & $\eu^{\hat{A}}\hat{B}\eu^{-\hat{A}} = \sum_{n=0}^\infty \frac{1}{n!}\mathrm{ad}_{\hat{A}}^n(\hat{B})$ & Lema de Hadamard (transformaciones unitarias de semejanza). \\
        \addlinespace[2pt]
        & $\eu^{\hat{A}}\eu^{\hat{B}} = \eu^{\hat{A}+\hat{B}}\eu^{\frac{1}{2}[\hat{A},\hat{B}]}$ & Fórmula BCH (válida si $[[\hat{A},\hat{B}],\hat{A}]=[[\hat{A},\hat{B}],\hat{B}]=0$). \\
        \midrule
        \textbf{Momento Angular} & $[\hat{x}_i, \hat{p}_j] = \iu\hbar\delta_{ij}\hat{I}$ & Relaciones canónicas cartesianas de Heisenberg. \\
        \addlinespace[2pt]
        & $\hat{L}_k = \varepsilon_{kij}\hat{x}_i\hat{p}_j$ & Definición tensorial del momento angular orbital. \\
        \addlinespace[2pt]
        & $[\hat{L}_i, \hat{L}_j] = \iu\hbar\,\varepsilon_{ijk}\hat{L}_k$ & Álgebra de Lie del grupo de rotaciones $\mathfrak{so}(3)$. \\
        \addlinespace[2pt]
        & $\hat{\vec{L}} \times \hat{\vec{L}} = \iu\hbar\,\hat{\vec{L}} \neq \vec{0}$ & Producto vectorial cuántico de autoacoplamiento. \\
        \addlinespace[2pt]
        & $[\hat{L}^2, \hat{L}_i] = 0, \quad \forall\,i$ & Invariancia del operador Casimir (magnitud total al cuadrado). \\
        \addlinespace[2pt]
        & $\hat{L}_\pm = \hat{L}_x \pm \iu\hat{L}_y$ & Operadores escalera de subida (+) y bajada (-). \\
        \addlinespace[2pt]
        & $[\hat{L}_z, \hat{L}_\pm] = \pm\hbar\hat{L}_\pm$ & Conmutadores fundamentales de escalera. \\
        \addlinespace[2pt]
        & $\hat{L}_\pm\ket{\ell, m} = \hbar\sqrt{\ell(\ell+1)-m(m\pm 1)}\ket{\ell, m\pm 1}$ & Acción espectral normalizada en la base $\{\ket{\ell, m}\}$. \\
        \midrule
        \textbf{Potencial Central} & $\nabla^2 = \frac{1}{r^2}\partial_r(r^2\partial_r) - \frac{\hat{L}^2}{\hbar^2 r^2}$ & Laplaciano esférico con operador Casimir angular. \\
        \addlinespace[2pt]
        & $\hat{H} = -\frac{\hbar^2}{2\mu}\frac{1}{r^2}\partial_r(r^2\partial_r) + \frac{\hat{L}^2}{2\mu r^2} + V(r)$ & Hamiltoniano central con barrera centrífuga cuántica. \\
        \midrule
        \textbf{Producto Tensorial} & $(\hat{A}_1 \otimes \hat{B}_1)(\hat{A}_2 \otimes \hat{B}_2) = (\hat{A}_1\hat{A}_2)\otimes(\hat{B}_1\hat{B}_2)$ & Composición de operadores en espacios producto. \\
        \addlinespace[2pt]
        & $\hat{\vec{J}} = \hat{\vec{J}}_1 + \hat{\vec{J}}_2, \quad [\hat{J}_i, \hat{J}_j] = \iu\hbar\,\varepsilon_{ijk}\hat{J}_k$ & Momento angular total acoplado. \\
        \addlinespace[2pt]
        & $|j_1 - j_2| \le j \le j_1 + j_2$ & Regla triangular de adición de momentos angulares. \\
        \midrule
        \textbf{CCOC} & $\{\hat{H}, \hat{L}^2, \hat{L}_z\}$ (orbital; con espín, $\{\hat{H}, \hat{L}^2, \hat{L}_z, \hat{S}^2, \hat{S}_z\}$) & Base ortonormal común $\ket{n, \ell, m_\ell}$ (orbital) o $\ket{n, \ell, m_\ell, s, m_s}$ en átomos hidrogenoides ($\hat{S}^2$ es constante trivial para $s=1/2$). \\
        \bottomrule
    \end{tabular}%
    }
\end{table}

\section{Respuestas y Guía Metodológica de los Ejercicios Propuestos}
\label{anexo:solucionario}

A continuación se presentan las resoluciones analíticas condensadas de los ejercicios propuestos de la Sección~\ref{sec:ejercicios} para corroborar los resultados obtenidos en el autoestudio.

\subsection[Solución 1: Conmutadores de exponenciales]{Solución al Ejercicio 1: Conmutadores de Exponenciales de Operadores}
\begin{enumerate}
    \item Para $n = 1$, $[\hat{A}, \hat{B}] = 1\hat{A}^0[\hat{A}, \hat{B}]$, lo cual es trivialmente cierto. Suponiendo válido para $n = k$, evaluamos para $k+1$:
    \begin{align}
        [\hat{A}^{k+1}, \hat{B}] &= [\hat{A}\hat{A}^k, \hat{B}] = \hat{A}[\hat{A}^k, \hat{B}] + [\hat{A}, \hat{B}]\hat{A}^k \nonumber\\
        &= \hat{A}(k\hat{A}^{k-1}[\hat{A}, \hat{B}]) + [\hat{A}, \hat{B}]\hat{A}^k = k\hat{A}^k[\hat{A}, \hat{B}] + \hat{A}^k[\hat{A}, \hat{B}] \nonumber\\
        &= (k+1)\hat{A}^k[\hat{A}, \hat{B}],
    \end{align}
    donde se utilizó que $[\hat{A}, \hat{B}]$ conmuta con $\hat{A}$. Por inducción, queda demostrado para todo $n \in \mathbb{N}$.
    \item Expandiendo la exponencial:
    \begin{align}
        [\eu^{\hat{A}}, \hat{B}] &= \sum_{n=0}^\infty \frac{1}{n!}[\hat{A}^n, \hat{B}] = \sum_{n=1}^\infty \frac{1}{n!} n\hat{A}^{n-1}[\hat{A}, \hat{B}] \nonumber\\
        &= [\hat{A}, \hat{B}]\sum_{n=1}^\infty \frac{\hat{A}^{n-1}}{(n-1)!} = [\hat{A}, \hat{B}]\sum_{m=0}^\infty \frac{\hat{A}^m}{m!} = [\hat{A}, \hat{B}]\eu^{\hat{A}}.
    \end{align}
\end{enumerate}

\subsection[Solución 2: Generador de dilataciones $\hat{D}$]{Solución al Ejercicio 2: Hermiticidad e Invariancia del Generador de Dilataciones $\hat{D}$}
\begin{enumerate}
    \item \textit{Hermiticidad y Conmutador con $\hat{L}_i$:} Tomando el adjunto hermítico de $\hat{D} = \frac{1}{2}(\hat{\vec{r}}\cdot\hat{\vec{p}} + \hat{\vec{p}}\cdot\hat{\vec{r}})$:
    \begin{equation}
        \hat{D}\adj = \frac{1}{2}\left( (\hat{\vec{r}}\cdot\hat{\vec{p}})\adj + (\hat{\vec{p}}\cdot\hat{\vec{r}})\adj \right) = \frac{1}{2}\left( \hat{\vec{p}}\cdot\hat{\vec{r}} + \hat{\vec{r}}\cdot\hat{\vec{p}} \right) = \hat{D}.
    \end{equation}
    Además, usando $[\hat{x}_k, \hat{p}_k] = \iu\hbar\hat{I}$ (sumando sobre $k=1,2,3$ da $3\iu\hbar\hat{I}$), se tiene $\hat{\vec{p}}\cdot\hat{\vec{r}} = \hat{\vec{r}}\cdot\hat{\vec{p}} - 3\iu\hbar\hat{I}$, lo que reescribe $\hat{D} = \hat{\vec{r}}\cdot\hat{\vec{p}} - \frac{3}{2}\iu\hbar\hat{I}$.\\
    Dado que el múltiplo escalar conmuta con todo operador, el conmutador con $\hat{L}_i$ se reduce a:
    \begin{equation}
        [\hat{L}_i, \hat{D}] = [\hat{L}_i, \hat{\vec{r}}\cdot\hat{\vec{p}}] = [\hat{L}_i, \hat{x}_k\hat{p}_k] = [\hat{L}_i, \hat{x}_k]\hat{p}_k + \hat{x}_k[\hat{L}_i, \hat{p}_k].
    \end{equation}
    Sustituyendo los resultados de operador vectorial $[\hat{L}_i, \hat{x}_k] = \iu\hbar\varepsilon_{ikm}\hat{x}_m$ y $[\hat{L}_i, \hat{p}_k] = \iu\hbar\varepsilon_{ikm}\hat{p}_m$:
    \begin{align}
        [\hat{L}_i, \hat{D}] &= (\iu\hbar\,\varepsilon_{ikm}\hat{x}_m)\hat{p}_k + \hat{x}_k(\iu\hbar\,\varepsilon_{ikm}\hat{p}_m) \nonumber\\
        &= \iu\hbar\,\varepsilon_{ikm}\left( \hat{x}_m\hat{p}_k + \hat{x}_k\hat{p}_m \right).
    \end{align}
    Como $\hat{x}_m$ y $\hat{p}_k$ conmutan cuando $m \neq k$, la cantidad entre paréntesis es estrictamente \textbf{simétrica} ante el intercambio de $m$ y $k$, mientras que $\varepsilon_{ikm}$ es estrictamente \textbf{antisimétrico}. Por el teorema de contracción simétrico-antisimétrico demostrado en la Sección~\ref{sec:momento_angular}, la contracción se anula idénticamente:
    \begin{equation}
        [\hat{L}_i, \hat{D}] = 0, \quad \forall\,i \in \{1, 2, 3\}.
    \end{equation}
    \item Como $\hat{L}^2 = \hat{L}_i\hat{L}_i$, aplicando Leibniz:
    \begin{equation}
        [\hat{L}^2, \hat{D}] = [\hat{L}_i\hat{L}_i, \hat{D}] = \hat{L}_i[\hat{L}_i, \hat{D}] + [\hat{L}_i, \hat{D}]\hat{L}_i = \hat{L}_i(0) + (0)\hat{L}_i = 0.
    \end{equation}
    Físicamente, $\hat{D}$ es el generador hermítico de dilataciones espaciales (transformaciones de escala $\vec{r} \to \eu^\lambda\vec{r}$); como las dilataciones son isotrópicas (escalares puros que no alteran los ángulos ni la dirección espacial), conmutan con todas las rotaciones del grupo $\mathrm{SO}(3)$.
\end{enumerate}

\subsection[Solución 3: Paquetes con chirp]{Solución al Ejercicio 3: Paquetes con Chirp y Saturación de Incertidumbre}
Para el estado $\psi(x) = C \exp\left(-\frac{x^2}{2\sigma^2} + \iu\frac{\alpha}{2}x^2\right)$:
\begin{enumerate}
    \item La densidad de probabilidad de posición es $|\psi(x)|^2 = |C|^2 \exp(-x^2/\sigma^2)$, idéntica al paquete sin fase. Por ende, $\ev{\hat{x}} = 0$ y $\ev{\hat{x}^2} = \sigma^2/2$, lo que arroja $\Delta x = \sigma/\sqrt{2}$.
    \item Al calcular $\hat{p}_x\psi(x) = -\iu\hbar\frac{\dif\psi}{\dif x} = -\iu\hbar\left(-\frac{x}{\sigma^2} + \iu\alpha x\right)\psi(x) = \left(\hbar\alpha x + \iu\frac{\hbar}{\sigma^2}x\right)\psi(x)$.\\
    Calculando la norma de este estado:
    \begin{equation}
        \ev{\hat{p}_x^2} = \left(\hbar^2\alpha^2 + \frac{\hbar^2}{\sigma^4}\right)\ev{\hat{x}^2} = \left(\hbar^2\alpha^2 + \frac{\hbar^2}{\sigma^4}\right)\frac{\sigma^2}{2} = \frac{\hbar^2}{2\sigma^2}(1 + \alpha^2\sigma^4).
    \end{equation}
    Por ende, $\Delta p_x = \frac{\hbar}{\sqrt{2}\sigma}\sqrt{1 + \alpha^2\sigma^4}$.
    \item La covarianza cuántica resulta ser:
    \begin{equation}
        \mathrm{Cov}(\hat{x}, \hat{p}_x) = \frac{1}{2}\ev{\hat{x}\hat{p}_x + \hat{p}_x\hat{x}} = \hbar\alpha\ev{\hat{x}^2} = \frac{\hbar\alpha\sigma^2}{2}.
    \end{equation}
    Evaluemos el miembro derecho de Robertson--Schrödinger:
    \begin{equation}
        \frac{1}{4}|\ev{\{\hat{x}, \hat{p}_x\}}|^2 + \frac{1}{4}|\ev{[\hat{x}, \hat{p}_x]}|^2 = \left(\frac{\hbar\alpha\sigma^2}{2}\right)^2 + \left(\frac{\hbar}{2}\right)^2 = \frac{\hbar^2}{4}(1 + \alpha^2\sigma^4).
    \end{equation}
    Evaluemos el producto de varianzas:
    \begin{equation}
        (\Delta x)^2 (\Delta p_x)^2 = \left(\frac{\sigma^2}{2}\right)\left(\frac{\hbar^2}{2\sigma^2}(1 + \alpha^2\sigma^4)\right) = \frac{\hbar^2}{4}(1 + \alpha^2\sigma^4).
    \end{equation}
    ¡Ambos lados son idénticos! El estado satura con igualdad matemática exacta la relación de Robertson--Schrödinger, demostrando que cuando existe covarianza no nula ($\alpha \neq 0$), la cota real de dispersión es estrictamente mayor que la cota básica de Heisenberg ($\hbar/2$).
\end{enumerate}

\subsection[Solución 4: Acoplamiento $LS$ y centro de gravedad]{Solución al Ejercicio 4: Acoplamiento Espín-Órbita y Regla del Centro de Gravedad}
\begin{enumerate}
    \item Expandiendo el cuadrado del operador momento angular total $\hat{\vec{J}} = \hat{\vec{L}} + \hat{\vec{S}}$:
    \begin{equation}
        \hat{J}^2 = (\hat{\vec{L}} + \hat{\vec{S}})\cdot(\hat{\vec{L}} + \hat{\vec{S}}) = \hat{L}^2 + \hat{S}^2 + \hat{\vec{L}}\cdot\hat{\vec{S}} + \hat{\vec{S}}\cdot\hat{\vec{L}}.
    \end{equation}
    Como los operadores espaciales conmutan con los operadores de espín ($[\hat{L}_i, \hat{S}_j] = 0$), el producto escalar conmuta: $\hat{\vec{S}}\cdot\hat{\vec{L}} = \hat{\vec{L}}\cdot\hat{\vec{S}}$. Por ende:
    \begin{equation}
        \hat{J}^2 = \hat{L}^2 + \hat{S}^2 + 2\hat{\vec{L}}\cdot\hat{\vec{S}} \implies \hat{\vec{L}}\cdot\hat{\vec{S}} = \frac{1}{2}\left( \hat{J}^2 - \hat{L}^2 - \hat{S}^2 \right).
    \end{equation}
    \item Para $\ell = 1$ y $s = 1/2$, la regla triangular establece $|1 - 1/2| \le j \le 1 + 1/2$, de donde:
    \begin{equation}
        j \in \{1/2, \, 3/2\}.
    \end{equation}
    Las etiquetas espectroscópicas $^{2s+1}L_j$ son:
    \begin{itemize}
        \item Para $j = 3/2$: $^{2}P_{3/2}$, con degeneración $g_{3/2} = 2(3/2) + 1 = 4$.
        \item Para $j = 1/2$: $^{2}P_{1/2}$, con degeneración $g_{1/2} = 2(1/2) + 1 = 2$.
    \end{itemize}
    Nótese que la dimensión total se conserva: $4 + 2 = 6 = (2\cdot 1 + 1)(2\cdot 1/2 + 1)$.
    \item Evaluando el valor propio de $\hat{H}_{SO} = \frac{\xi}{2}(\hat{J}^2 - \hat{L}^2 - \hat{S}^2)$ (donde $\xi \equiv \ev{\xi(r)}_{2p}$ es el valor esperado radial de la función de acoplamiento espín-órbita de la Sección~\ref{subsec:acoplamiento_espin_orbita}) sobre el estado acoplado $\ket{j, M}$ con $\ell=1$ y $s=1/2$:
    \begin{align}
        \Delta E(j) &= \frac{\xi\hbar^2}{2}\left[ j(j+1) - \ell(\ell+1) - s(s+1) \right] \nonumber\\
        &= \frac{\xi\hbar^2}{2}\left[ j(j+1) - 2 - \frac{3}{4} \right] = \frac{\xi\hbar^2}{2}\left[ j(j+1) - \frac{11}{4} \right].
    \end{align}
    \begin{itemize}
        \item Para $j = 3/2$: $j(j+1) = \frac{3}{2}\cdot\frac{5}{2} = \frac{15}{4}$, luego:
        \begin{equation}
            \Delta E(^2P_{3/2}) = \frac{\xi\hbar^2}{2}\left( \frac{15}{4} - \frac{11}{4} \right) = \frac{\xi\hbar^2}{2}(1) = +\frac{1}{2}\xi\hbar^2.
        \end{equation}
        \item Para $j = 1/2$: $j(j+1) = \frac{1}{2}\cdot\frac{3}{2} = \frac{3}{4}$, luego:
        \begin{equation}
            \Delta E(^2P_{1/2}) = \frac{\xi\hbar^2}{2}\left( \frac{3}{4} - \frac{11}{4} \right) = \frac{\xi\hbar^2}{2}(-2) = -\xi\hbar^2.
        \end{equation}
    \end{itemize}
    El desdoblamiento energético total resulta ser:
    \begin{equation}
        \Delta E_{\text{fina}} = \Delta E(^2P_{3/2}) - \Delta E(^2P_{1/2}) = \frac{1}{2}\xi\hbar^2 - (-\xi\hbar^2) = \frac{3}{2}\xi\hbar^2.
    \end{equation}
    \item Evaluando la suma ponderada por las degeneraciones:
    \begin{equation}
        \sum_j (2j+1)\Delta E(j) = 4\left(+\frac{1}{2}\xi\hbar^2\right) + 2(-\xi\hbar^2) = +2\xi\hbar^2 - 2\xi\hbar^2 = 0.
    \end{equation}
    Físicamente, la perturbación interna de espín-órbita redistribuye los niveles de energía pero no altera la energía media (la traza del operador en el subespacio de dimensión 6 es idénticamente nula).

\end{enumerate}

\subsection[Solución 5: Estados de Bell y traza parcial]{Solución al Ejercicio 5: Producto Tensorial, Estados de Bell y Traza Parcial}
\begin{enumerate}
    \item \textit{Demostración de entrelazamiento por contradicción:}
    Supongamos que $\ket{\Phi^+}$ es separable, es decir, existen $\ket{\psi_A} = a_0\ket{0} + a_1\ket{1}$ y $\ket{\psi_B} = b_0\ket{0} + b_1\ket{1}$ tales que:
    \begin{equation}
        \ket{\psi_A} \otimes \ket{\psi_B} = a_0 b_0\ket{00} + a_0 b_1\ket{01} + a_1 b_0\ket{10} + a_1 b_1\ket{11} = \frac{1}{\sqrt{2}}\ket{00} + \frac{1}{\sqrt{2}}\ket{11}.
    \end{equation}
    Igualando coeficientes base a base:
    \begin{equation}
        a_0 b_0 = \frac{1}{\sqrt{2}}, \qquad a_0 b_1 = 0, \qquad a_1 b_0 = 0, \qquad a_1 b_1 = \frac{1}{\sqrt{2}}.
    \end{equation}
    De $a_0 b_0 = 1/\sqrt{2} \neq 0$, se deduce que $a_0 \neq 0$ y $b_0 \neq 0$. De $a_0 b_1 = 0$ con $a_0 \neq 0$, se tiene forzosamente $b_1 = 0$. Pero si $b_1 = 0$, entonces $a_1 b_1 = 0$, lo cual contradice directamente que $a_1 b_1 = 1/\sqrt{2} \neq 0$. La contradicción demuestra que $\ket{\Phi^+}$ no puede escribirse como un producto tensorial simple; por ende, es un \textbf{estado entrelazado}.
    \item \textit{Matriz de densidad global de $4\times 4$:}
    En la base ordenada $\{\ket{00}, \ket{01}, \ket{10}, \ket{11}\}$, el vector de estado es $\ket{\Phi^+} = \frac{1}{\sqrt{2}}\begin{pmatrix} 1 & 0 & 0 & 1 \end{pmatrix}^T$. Por consiguiente:
    \begin{equation}
        \hat{\rho} = \ket{\Phi^+}\bra{\Phi^+} = \frac{1}{2}\begin{pmatrix} 1 \\ 0 \\ 0 \\ 1 \end{pmatrix}\begin{pmatrix} 1 & 0 & 0 & 1 \end{pmatrix} = \frac{1}{2}\begin{pmatrix}
            1 & 0 & 0 & 1 \\
            0 & 0 & 0 & 0 \\
            0 & 0 & 0 & 0 \\
            1 & 0 & 0 & 1
        \end{pmatrix}.
    \end{equation}
    \item \textit{Traza parcial y matriz densidad reducida:}
    Evaluamos la traza parcial sobre el subsistema $B$:
    \begin{align}
        \hat{\rho}_A &= \sum_{k=0}^1 \bra{k}_B \hat{\rho} \ket{k}_B = \bra{0}_B \hat{\rho} \ket{0}_B + \bra{1}_B \hat{\rho} \ket{1}_B \nonumber\\
        &= \frac{1}{2}\ket{0}_A\bra{0}_A + \frac{1}{2}\ket{1}_A\bra{1}_A = \frac{1}{2}\begin{pmatrix} 1 & 0 \\ 0 & 1 \end{pmatrix} = \frac{1}{2}\hat{I}_2.
    \end{align}
    Como $\hat{\rho}_A = \frac{1}{2}\hat{I}_2$, se tiene $\mathrm{Tr}(\hat{\rho}_A^2) = \mathrm{Tr}\left(\frac{1}{4}\hat{I}_2\right) = \frac{2}{4} = \frac{1}{2} < 1$. El subsistema local $A$ se encuentra en un \textbf{estado maximalmente mezclado}: la probabilidad de medir $\ket{0}$ o $\ket{1}$ es exactamente $1/2$ (entropía de von Neumann máxima $S = -\mathrm{Tr}(\rho_A \log_2 \rho_A) = 1$ bit). Toda la certeza del estado global está contenida en el entrelazamiento cuántico no local.
\end{enumerate}

\end{appendices}
